% 课程论文 LaTeX 模板（带注释版）
% 使用方法：
% 1. 复制本文件，改成你的论文文件名。
% 2. 修改下面的标题、作者信息、摘要、关键词和正文。
% 3. 建议使用 XeLaTeX 编译，因为中文支持最稳定。

% ctexart 是中文论文/文章常用文档类。
% UTF8 表示文件编码为 UTF-8；zihao=-4 表示正文小四号；a4paper 表示 A4 纸。
\documentclass[UTF8,zihao=-4,a4paper]{ctexart}

% geometry：设置页面边距。
% setspace：设置行距。
% indentfirst：让每个章节开头第一段也首行缩进。
% graphicx：插入图片时使用。
% booktabs：制作三线表时使用。
% amsmath：数学公式支持。
% hyperref：生成 PDF 链接和目录跳转信息。
\usepackage{geometry}
\usepackage{setspace}
\usepackage{indentfirst}
\usepackage{graphicx}
\usepackage{booktabs}
\usepackage{amsmath}
\usepackage{fontspec}
\usepackage{hyperref}

% 字体设置：建议使用 XeLaTeX 编译。
% 按课程论文常用要求设置：中文正文宋体，英文与数字 Times New Roman，标题黑体。
\setmainfont{Times New Roman}
\setsansfont{Times New Roman}
\setmonofont{Times New Roman}
\setCJKmainfont{SimSun}
\setCJKsansfont{SimHei}
\setCJKmonofont{FangSong}
\setCJKfamilyfont{zhsong}{SimSun}
\setCJKfamilyfont{zhhei}{SimHei}
\setCJKfamilyfont{zhkai}{KaiTi}
\setCJKfamilyfont{zhfs}{FangSong}
\newCJKfontfamily\titlehei{SimHei}

% 页面边距：上、下 2.54cm，左、右 3.17cm。
\geometry{left=3.17cm,right=3.17cm,top=2.54cm,bottom=2.54cm}

% 行距设置为 1.5 倍。
\setstretch{1.5}

% 段落首行缩进 2 个中文字符。
\setlength{\parindent}{2em}

% 摘要样式：左对齐，标签黑体小四，摘要正文宋体小四。
\renewenvironment{abstract}
  {\par\noindent{\titlehei\zihao{-4}摘要：}\songti\zihao{-4}}
  {\par\vspace{0.5em}}

% PDF 链接设置。
% colorlinks=true 表示链接使用颜色而不是方框。
% linkcolor/citecolor/urlcolor 分别控制内部链接、引用链接和网址颜色。
\hypersetup{
  colorlinks=true,
  linkcolor=black,
  citecolor=black,
  urlcolor=blue,
  pdfauthor={zpern},
  pdftitle={title}
}

% 设置章节标题样式。
% section 对应 \section{...}，也就是一级标题。
% number = \chinese{section} 表示章节编号用“一、二、三”。
% name = {,、} 表示标题显示为“一、引言”这种格式。
\ctexset{
  section = {
    format = \zihao{4}\titlehei,
    name = {,、},
    number = \chinese{section}
  },
  subsection = {
    format = \zihao{-4}\titlehei
  }
}

% 论文标题。
% \\ 表示换行；主标题三号黑体，副标题稍微小一点。
\title{\titlehei\zihao{3} 这里填写论文主标题\\
\large 这里填写副标题，没有副标题可以删除这一行}

% 作者信息。
% \underline{\makebox[3cm][c]{...}} 会生成固定宽度的下划线。
% 如果是小组作业，可以写“小组：第X组”；如果是个人作业，可以写姓名、学号、学院。
\author{
姓名：\underline{\makebox[3cm][c]{张三}}
\quad
学号：\underline{\makebox[3cm][c]{12345678}}\\
学院：\underline{\makebox[3cm][c]{某某学院}}
}

% 日期。
% 可以写固定日期，如 2026年5月；也可以写 \today 自动显示当天日期。
\date{2026年5月}

% 正文从这里开始。
\begin{document}

% 生成标题、作者、日期。
\maketitle

% 摘要环境。
% 摘要一般写研究对象、研究思路、主要观点和结论，课程论文可写 150-300 字左右。
\begin{abstract}
  这里填写摘要内容。摘要应当简要说明本文研究什么问题、依据什么理论或材料展开分析、
  得出了什么主要观点。摘要不宜写得太口语化，也不要展开太多细节。
\end{abstract}

% 关键词。
% \noindent 表示这一行不缩进；关键词标签使用黑体小四。
% 关键词之间通常用中文分号隔开。
\noindent{\titlehei\zihao{-4}关键词：}{\songti\zihao{-4}关键词一；关键词二；关键词三；关键词四}

% 一级标题使用 \section{标题名}。
\section{引言}

这里写引言。引言一般说明选题背景、问题意识、研究意义和文章结构。
如果是课程论文，可以简要联系课程主题，说明为什么这个问题值得讨论。

\section{理论基础或问题分析}

这里写正文第一部分。每一段之间空一行即可，LaTeX 会自动处理段落缩进。
不要手动敲很多空格来排版。

% 二级标题使用 \subsection{标题名}。
\subsection{二级标题示例}

这里写二级标题下的内容。一般只有当某一节内容较多时，才需要使用二级标题。

\section{案例分析或实践展开}

这里写正文第二部分。可以结合具体案例、政策材料、课程汇报内容或现实问题展开分析。

% 插入图片示例。
% 如果暂时不用图片，可以保留这段注释，不需要删除。
% 使用图片时，把图片放到同一目录或 figures 文件夹中，然后取消下面几行前面的 %。
%
% \begin{figure}[htbp]
%   \centering
%   \includegraphics[width=0.7\textwidth]{figures/example.png}
%   \caption{这里填写图片标题}
% \end{figure}

% 插入表格示例。
% booktabs 提供 \toprule、\midrule、\bottomrule，用来做更规范的三线表。
%
% \begin{table}[htbp]
%   \centering
%   \caption{这里填写表格标题}
%   \begin{tabular}{lll}
%     \toprule
%     指标 & 内容 & 说明 \\
%     \midrule
%     示例一 & 示例内容 & 示例说明 \\
%     示例二 & 示例内容 & 示例说明 \\
%     \bottomrule
%   \end{tabular}
% \end{table}

\section{结论}

这里写结论。结论一般概括全文观点，回应引言中提出的问题，
也可以简要指出现实启示或后续需要继续研究的问题。

% 参考文献。
% thebibliography 后面的 9 表示文献编号宽度，文献少于 10 条时用 9 即可。
% \bibitem{ref1} 中的 ref1 是文献标签，可以自己命名。
\begin{thebibliography}{9}
  \bibitem{ref1}
  作者：《书名或文章名》，出版社或网站，年份。

  \bibitem{ref2}
  作者：《文章名》，《期刊名》，年份，第几期。

  \bibitem{ref3}
  网站或机构：《资料标题》，发布日期或访问日期。
\end{thebibliography}

% 正文到这里结束。
\end{document}
